\section{共享功率平方上界：表示压缩与两条历史诊断}
\label{sec:shared-power-square-upper}

共享功率平方图在两条固定历史上显著缩小上界主问题，均获得可用的最大化全局求解界（best bound）。
H1 的动作代理下界略降、H2 略升；两条历史的共享动作相对恒充动作增益区间都跨过零。
因此，本阶段支持上界表示的数值构造与计算效果，动作辨别收益排序仍未确定。
既有 C 优化诊断中，515 次原 \texttt{secant\_upper} 全部超时，均缺少 incumbent 和 best bound；
其全域代理上界保留初始 control-box 见证上界。
有限方向下界主问题的最优性只针对其有限方向代理，不能据此称旧方法在完整平方代理域上全局最优。

\subsection{同协议的两条 H-only 搜索}
本阶段只使用经白名单检查的历史 $P,Q,V$、总表 $P,Q$ 及根电压六个通道。
两条历史分别以新种子 28101301/28101302 构造四候选、每候选两中心、24 个冻结 pair；
中心构造使用 $\beta$ 的 $\pm30\%$ 局部域、同点源调整 $\pm0.05$、均值预测，
并允许保留未优化但实际可行的低损失种子。
这些有限中心属于设计参数域，不能解释为置信样本或完整 MLE 域。
原版与共享版在每条历史中使用相同 pair 指纹、8 时段单 EV 的 6 kWh 运行模型，
共同初始化仅为站点 0 的 3 kW 恒充及零动作。
搜索预算固定为最多 3 轮、每次 MILP 3 秒、总预算 25 秒、初始 2 段最多 8 段，
停止尺度为绝对 $10^{-4}$ 加相对 $10^{-3}$。

适配器复制原 \texttt{adaptive\_search} 的函数全局表，保留同一 code object、默认参数与关键字默认参数；
双方共同使用透传 QP 记录器，策略差异仅为 \texttt{upper\_master} 后端。
原全局函数和求解器在每次调用后恢复，15 个冻结核心源码哈希前后一致。
后续上界动作不同会产生不同 cuts 与见证池，因而跨轮模型规模是完整搜索的描述。
此前 H1 同池单轮比较另固定 24 pair、每 pair 一个恒充见证及两段 PWL，
得到 $5914/4932$（变量/行）至 $322/276$ 的压缩；原版超时缺界，共享版数值全局求解界为 52.9025575268。

\begin{table}[htbp]
\centering\small
\caption{两条固定历史的完整搜索结果。上下界仅对应冻结仿射有限 pair 代理。}
\label{tab:shared-upper-search}
\begin{tabular}{llrrrrl}
\hline
历史 & 后端 & 动作 LB & 全域 UB & 原始 gap & 秒 & 上界变量数 \\
\hline
710009 & 原版 & 52.902544 & 52.902663 & 0.000119 & 10.7696 & 23194, 57754 \\
710009 & 共享 & 52.889957 & 52.902573 & 0.012616 & 4.1538 & 322, 466 \\
720001 & 原版 & 51.107621 & 51.149059 & 0.041437 & 6.9915 & 17434 \\
720001 & 共享 & 51.133666 & 51.148963 & 0.015297 & 2.1541 & 322 \\
\hline
\end{tabular}
\end{table}

表中只保留六位小数，原始数值和完整求解界保存在归档。
本 pilot 的三次原 upper 均超时且缺少 incumbent/best bound，三次共享 upper 均返回二者。
共享 H1 的最后一轮 status=0 仍记录非零 \texttt{mip\_gap}=$4.1239\times10^{-7}$，保留该原值。
四次搜索均按 \texttt{proxy\_gap\_met} 停止，但 raw gap 都超过 $10^{-4}$；
停止由 $10^{-4}+10^{-3}\max(1,|\mathrm{LB}|)$ 容许，分母包含历史代价。
所选动作活动 AC 见证的未来平方距离占总代价比例，H1 原版/共享约为
$2.28\times10^{-7}/2.06\times10^{-7}$，H2 约为 $4.13\times10^{-6}/6.09\times10^{-6}$。
该占比描述保存的活动见证，不能扩展为统一非线性 infimum 的分解。

\subsection{时段支撑与单 EV 排除交叉项}
令 $p$ 表示冻结 pair，$z\in\mathcal Z_p$ 表示可行竞争调整变量，$u=(u_{ti})_{t,i}$ 为时段优先排列的站点净功率。
固定 $z$ 时的目标为
\begin{equation}
 f_{p,z}(u)=\|c_p+A_pz+B_pu\|_2^2+\|F_pz\|_2^2,
 \qquad Q(u)=\min_p\inf_{z\in\mathcal Z_p}f_{p,z}(u).
\end{equation}
记 $r_{p,z}=c_p+A_pz$，$b_{p,ti}=B_p[:,(t,i)]$。
实际 $B_p$ 的每个残差行精确非零支撑至多属于一个时段，故不同时间的列满足
$b_{p,ti}^{\mathsf T}b_{p,sj}=0$（$t\ne s$）。
单 EV 单位路径及停留功率绑定保证每时段至多一个站点净功率非零，即
$u_{ti}u_{tj}=0$（$i\ne j$）。
展开得到
\begin{align}
 \|B_pu\|_2^2
 &=\sum_{t,i}\|b_{p,ti}\|_2^2u_{ti}^2
   +2\!\sum_{(t,i)<(s,j)}b_{p,ti}^{\mathsf T}b_{p,sj}u_{ti}u_{sj}\\
 &=\sum_{t,i}d_{p,ti}u_{ti}^2,\qquad d_{p,ti}=\|b_{p,ti}\|_2^2\ge0,\\
 f_{p,z}(u)&=a_{p,z}+\ell_{p,z}^{\mathsf T}u+\sum_{t,i}d_{p,ti}u_{ti}^2,\\
 a_{p,z}&=\|r_{p,z}\|_2^2+\|F_pz\|_2^2,\\
 \ell_{p,z}&=2B_p^{\mathsf T}r_{p,z}.
\end{align}
同期 $B_p$ 列可以不正交；交叉项消失依赖运行域的功率互斥。
实现对所有实际 $B_p$ 逐行检查精确支撑，任何跨时段非零（包括 $10^{-300}$）均拒绝压缩，
不删除 near-zero 系数。多 EV 同期激活多个站点时，上式一般失效。

\subsection{共享 PWL 平方图与上下界方向}
对坐标 $j=(t,i)$ 的功率盒 $[l_j,h_j]$，取包含零的节点；
在相邻节点 $[a,b]$ 上定义弦线平方图
\begin{equation}
 s_j(u_j)=(a+b)u_j-ab,
 \qquad s_j(u_j)-u_j^2=(u_j-a)(b-u_j)\in[0,(b-a)^2/4].
\end{equation}
全局图以分段二进制及相邻端点权重精确编码。
零节点保证不活动坐标 $u_j=0$ 时 $s_j=0$。
所有 pair 和有限可行见证库 $Z_p\subset\mathcal Z_p$ 共享一套 $s_j$，每个见证只增加线性上界行
\begin{equation}
 T\le a_{p,z}+\ell_{p,z}^{\mathsf T}u+\sum_jd_{p,j}s_j(u_j).
\end{equation}
令 $W(u)=\min_{p,z\in Z_p}f_{p,z}(u)$，$S(u)$ 为相应弦线最小值，$\Delta_j$ 为坐标最大段宽。
则在单 EV 可行域内
\begin{equation}
 Q(u)\le W(u)\le S(u),\qquad
 0\le S(u)-W(u)\le\max_p E_p,
 \quad E_p=\sum_t\max_i\frac{d_{p,ti}\Delta_{ti}^2}{4}.
\end{equation}
证明中，$Q\le W$ 来自有限见证缩小最小化域；$W\le S$ 来自 $d_{p,j}\ge0$ 的弦线误差方向。
取达到 $W(u)$ 的有限见证可得最后一个误差界。
该 $E_p$ 只控制 PWL 平方误差，有限见证误差 $W-Q$ 需另行处理。
最大化该上界主问题的全局求解界（乘回 scale）给出 $\max_{u\in\mathcal U}Q(u)$ 的数值上界；
仅有主问题 incumbent 时不能把其当作全域上界。
另外的 $T\le\mathrm{qbar}$ 使用原 control-box 见证上界，仍保持对 $Q$ 的上界方向。
浮点 RHS guard 属数值措施，本证据没有区间算术证书。
每坐标使用 $m$ 段时，共享图新增变量为 $J(3m+1)$，另有一个目标变量，
其中 $J$ 为功率坐标数；变量计数与 pair/见证/残差行数无关，见证约束数仍随见证池增长。

\subsection{动作区间、AC 回代及局限}
固定动作的 QP 可行 primal 给上界 $q_U$，盒约束切线给数值下界 $q_L$。
相对恒充动作 $u_0$ 的有效代理增益区间为
$[q_L(u)-q_U(u_0),\ q_U(u)-q_L(u_0)]$。
\begin{table}[htbp]
\centering\small
\caption{所选动作相对恒充的代理增益区间。}
\label{tab:shared-upper-gain}
\begin{tabular}{llrr}
\hline
历史 & 后端 & 增益下端 & 增益上端 \\
\hline
710009 & 原版 & 0.00001137 & 0.10868240 \\
710009 & 共享 & $-0.01257494$ & 0.10868335 \\
720001 & 原版 & $-0.04127440$ & 0.15875763 \\
720001 & 共享 & $-0.01522992$ & 0.15877361 \\
\hline
\end{tabular}
\end{table}
共享所选动作减原版所选动作的区间分别为
$[-0.0125880823,\ 0.0000027350]$ 和 $[-0.0152442295,\ 0.0413047005]$，均跨零。
原版 H1 的小正下端 $1.14\times10^{-5}$ 应保留，但其量级及代理范围不足以支持实际未来辨别改善。
四个所选动作的最大逐 pair primal--tangent gap 依次约为
$1.78\times10^{-6}$、$0.0125890$、$0.0438844$、$0.0162590$；
上界主问题缩小并不能自动闭合竞争调整 QP 区间。

四个动作的 96 个 pair 均完成 H-only 完整 AC 回代，均物理可行。
可行回代值只给固定动作、固定 pair 非线性竞争调整最小化的见证上界；
实际未来排错、拓扑正确概率及完整 AC 设计全局界均未评价。
活动见证虽有较小 Taylor 误差，全部 pair 的最大 Taylor 范数仍约为
28.28、51.60、24.46、24.85，显示局部仿射代理缺少全域非线性误差控制。
原始功率保留 $-0.0$ 及共享 H1 的 $6.000000000000001$，未裁剪，
运行可行性按既定 $2\times10^{-6}$ 容差核查。

独立包内核验器只读取归档矩阵、QP $z$、MILP $x$ 与保存的 H 拟合 AC 状态；
不调用拟合优化器或 selector。
两 H 的 696 条 QP primal/tangent、16 个 MILP snapshot 和 96 个 AC 见证全部重放通过：
primal/tangent 最大差为零，snapshot 最大行违约 $1.18\times10^{-9}$，
AC cost 最大差 $9.09\times10^{-13}$，Taylor 范数最大差 $7.11\times10^{-15}$。
同池单轮另核验 72 条 QP 与两个 snapshot。
归档保留原字节 SHA256、大型 JSON 的无损 gzip 及解压后 SHA256；
源码快照与原核验器可读，另提供使用正确包导入的可迁移适配器和核验入口。
本次是两条固定 H 的计算诊断，运行顺序及单机墙钟限制其耗时外推；
原 C/D 策略、门限和统计结果保持冻结，后续正式统计需要单独冻结新协议。
